\chapter{Installation}
\label{ch:install}

\index{installation}

Acme builds from source on Linux, macOS, and Windows.  The build
produces a single binary, \file{bin/acme}, with no runtime
configuration files required.

\section{Linux}
\label{sec:install-linux}

\index{installation!Linux}

\subsection{Prerequisites}

You need a C compiler (GCC or Clang), GNU make, and the development
libraries for SDL3, FreeType, fontconfig, and zlib.

\paragraph{Debian / Ubuntu:}
\begin{lstlisting}[language=bash]
sudo apt install gcc make
sudo apt install libsdl3-dev libfontconfig1-dev \
                 libfreetype-dev zlib1g-dev
\end{lstlisting}

\paragraph{Fedora:}
\begin{lstlisting}[language=bash]
sudo dnf install gcc make
sudo dnf install SDL3-devel fontconfig-devel \
                 freetype-devel zlib-devel
\end{lstlisting}

\paragraph{Arch Linux:}
\begin{lstlisting}[language=bash]
sudo pacman -S gcc make sdl3 fontconfig freetype2 zlib
\end{lstlisting}

Note: SDL3 is relatively new.  If your distribution does not yet
package it, you may need to build SDL3 from source or use a PPA.

\subsection{Building}

\begin{lstlisting}[language=bash]
git clone <repository-url> acme
cd acme
make
\end{lstlisting}

The build compiles 18 static libraries and the acme binary.  The result
is \file{bin/acme}.

\subsection{Running}

\begin{lstlisting}[language=bash]
./bin/acme
\end{lstlisting}

You can pass files to open on the command line:

\begin{lstlisting}[language=bash]
./bin/acme main.c util.c README.md
\end{lstlisting}

To install system-wide, copy the binary:

\begin{lstlisting}[language=bash]
sudo cp bin/acme /usr/local/bin/
\end{lstlisting}

\section{macOS}
\label{sec:install-macos}

\index{installation!macOS}

\subsection{Prerequisites}

Install Xcode Command Line Tools and Homebrew packages:

\begin{lstlisting}[language=bash]
xcode-select --install
brew install sdl3 fontconfig freetype pkg-config
\end{lstlisting}

\subsection{Building and Running}

\begin{lstlisting}[language=bash]
make
./bin/acme
\end{lstlisting}

The build uses \cmd{pkg-config} to locate Homebrew libraries.  If
\cmd{pkg-config} is not available, it falls back to default Homebrew
paths (\file{/opt/homebrew} on Apple Silicon, \file{/usr/local} on
Intel Macs).

\subsection{Mouse on macOS}

\index{mouse!macOS}
Most Mac users do not have a three-button mouse.  Acme provides
modifier-click emulation:

\begin{itemize}
\item \key{Ctrl}-click = \button{2} (middle button / execute)
\item \key{Alt/Option}-click = \button{3} (right button / look)
\end{itemize}

A standard USB or Bluetooth three-button mouse is strongly recommended
for extended use.  See Chapter~\ref{ch:mouse} for details.

\section{Windows}
\label{sec:install-windows}

\index{installation!Windows}

Acme on Windows can be built two ways:

\begin{itemize}
\item \textbf{MSYS2 build} — uses the MSYS2 POSIX compatibility layer.
  Easier to build but requires MSYS2 installed at runtime.
\item \textbf{Native standalone build} — produces a single
  \file{acme.exe} that runs on any Windows~10 or later machine with
  no other files required.
\end{itemize}

Both builds require MSYS2 at \emph{build} time for the toolchain.

\subsection{Prerequisites}

\begin{enumerate}
\item Install MSYS2 from \url{https://www.msys2.org/}.
\item Open any MSYS2 shell from the Start menu (MSYS, MINGW64, or
      UCRT64 — the build system pins the right compilers automatically).
\item Install dependencies:
\end{enumerate}

\begin{lstlisting}[language=bash]
pacman -Syu
pacman -S git gcc make pkg-config zlib-devel
pacman -S mingw-w64-x86_64-gcc \
          mingw-w64-x86_64-sdl3 \
          mingw-w64-x86_64-fontconfig \
          mingw-w64-x86_64-freetype
\end{lstlisting}

\subsection{Building (MSYS2 build)}

From any MSYS2 shell:

\begin{lstlisting}[language=bash]
export PKG_CONFIG_PATH=/mingw64/lib/pkgconfig
export PKG_CONFIG_ALLOW_SYSTEM_CFLAGS=1
export PKG_CONFIG_ALLOW_SYSTEM_LIBS=1
make
./bin/acme
\end{lstlisting}

The binary depends on the MSYS2 runtime (\file{msys-2.0.dll}) and
mingw64 DLLs.  It must be run from an MSYS2 environment or with
those DLLs in the PATH.

See \file{Windows-MSYS.md} in the source tree for full details.

\subsection{Building (native standalone)}

From any MSYS2 shell:

\begin{lstlisting}[language=bash]
make native-win
\end{lstlisting}

This produces \file{bin/acme.exe} (approximately 13~MB), statically
linked with no MSYS2 or Cygwin runtime dependency.  Copy
\file{acme.exe} to any Windows~10 or later machine and run it
directly --- no installation, no other files required.

The default font is Consolas, included with Windows.  Fonts are
discovered automatically from \file{C:\textbackslash Windows\textbackslash Fonts}.
The \cmd{Win} command opens an interactive \cmd{cmd.exe} session
inside an acme window using the Windows ConPTY API.

See \file{Windows-native.md} for full details.

\subsection{Mouse on Windows}

The same modifier-click emulation is available:

\begin{itemize}
\item \key{Ctrl}-click = \button{2} (middle button)
\item \key{Alt}-click = \button{3} (right button)
\end{itemize}

Any standard three-button USB mouse works natively.

\section{Command-Line Options}
\label{sec:cmdline}

\index{command-line options}

\begin{longtable}{@{}lp{3.5in}@{}}
\toprule
Flag & Description \\
\midrule
\endhead
\cmd{-a} & Enable auto-indent for all windows at startup. \\
\cmd{-c \textit{N}} & Start with \textit{N} columns (default: 2). \\
\cmd{-f \textit{font}} & Set the proportional (variable-width) font. \\
\cmd{-F \textit{font}} & Set the fixed-width font. \\
\cmd{-l \textit{file}} & Load session state from a dump file. \\
\cmd{-r} & Reverse scrollbar button mapping (Button~1 scrolls down). \\
\cmd{-W \textit{WxH}} & Set initial window size, e.g., \cmd{-W 1200x800}. \\
\bottomrule
\end{longtable}

\noindent
Fonts are specified as fontconfig names with a size suffix.  The
default for both \cmd{-f} and \cmd{-F} is
\file{DejaVuSansMono/13a/font}.

\section{Environment Variables}
\label{sec:envvars}

\index{environment variables}

\begin{longtable}{@{}lp{3.3in}@{}}
\toprule
Variable & Description \\
\midrule
\endhead
\env{HOME} & User home directory.  Used to find \file{\textasciitilde/lib/plumbing}. \\
\env{acmeshell} & Shell used to run commands.  Defaults to \cmd{sh}. \\
\env{tabstop} & Default tab width in columns.  Default is 4. \\
\env{SHELL} & Shell used by the Win command. \\
\env{NAMESPACE} & Plan~9 namespace directory.  Usually auto-detected. \\
\bottomrule
\end{longtable}

\noindent
When acme runs an external command, it sets several variables in the
child's environment:

\begin{longtable}{@{}lp{3.3in}@{}}
\toprule
Variable & Description \\
\midrule
\endhead
\env{winid} & Numeric ID of the current window. \\
\env{\%} & Filename of the current window. \\
\env{samfile} & Same as \env{\%}. \\
\bottomrule
\end{longtable}
